Временная и пространственная сложность алгоритма нужны для оценки и сравнения программ между собой. 

\Def Временная сложность - величина, характеризующая количество элементарных действий программы, то есть тех, которые работают условно за 1 такт работы процессора. 

Модель памяти(RAM-модель):
\begin{enumerate}
    \item Память вычислителя безгранична.
    \item Доступ к памяти на чтение/запись требует пренебрежимо малого времени(округляется единичкой)
\end{enumerate}

\Def Пространственная сложность - величина, характеризующая количество единиц доп.памяти, которые необходимы для выполнения программы.

\subsection*{Асимптотические сравнения}

\Def Пусть имеются две функции $f(n)$  и $g(n)$, при этом $f, g \; : \N \rightarrow \N$,
тогда считается, что $f(n) = \mathcal{O} (g(n))$, если $\exists C > 0 \; \ \exists N_{0}: \ \forall n > N_{0} \; \; f(n) \leqslant C \cdot g(n)$

\Def Пусть имеются две функции $f(n)$  и $g(n)$, при этом $f, g \; : \N \rightarrow \N$,
тогда считается, что $f(n) = \Omega (g(n))$, если $\exists C > 0 \; \ \exists N_{0}: \ \forall n > N_{0} \; \; f(n) \geqslant C \cdot g(n)$

\Def Пусть имеются две функции $f(n)$  и $g(n)$, при этом $f, g \; : \N \rightarrow \N$,
тогда считается, что $f(n) = \Theta (g(n))$, если $\exists C_{1}, C_{2} > 0 \; \ \exists N_{0}: \ \forall n > N_{0} \; \; C_{1} \cdot g(n) \leqslant f(n) \leqslant C_{2} \cdot g(n)$

\subsection*{Мастер-теорема}

\Thbd Пусть имеется рекуррента:

$$\begin{cases}
aT(\frac{n}{b}) + O(n^c)\text{,} & n > 1\\
O(1), & n = 1
\end{cases}$$\\
где $a \in \N,\,b \in \R,\,b > 1,\,c \in \R^{+}$
Тогда:
\begin{itemize}
    \item Если $c > log_b\,a$, то $T(n) = O(n^c)$
    \item Если $c = log_b\,a$, то $T(n) = O(n^c\,log\,n)$
    \item Если $c < log_b\,a$, то $T(n) = O(n ^ {log_b\,a})$
\end{itemize}

\pagebreak